\section{Summary on One Selected Paper}

\subsection{Deep Robust Reinforcement Learning for Practical Algorithmic Trading}
\textit{Yang Li, Wanshan Zheng and Zibin Zheng (2019)}

In their 2019 paper, Yang Li, Wanshan Zheng and Zibin Zheng present a further development
of previously used RL algorithms such as A3C and DQN for application in stock trading, a
so-called \emph{deep robust reinforcement learning} framework. To cope with the problem of
uncontrollable external influences, called \emph{noise} in their study, the following steps
have to be performed first:
\begin{itemize}
  \item Sampling random lengths of time units.
  \item To limit the influence of news, which is usually highest at opening time, the time
        points between 9:30 am to 11:30 am and 13:00 pm to 15:00 pm are omitted.
  \item Removal of the low volatility from the calculation in order to mitigate the effect
        caused by individual investors (as opposed to institutional investors).
\end{itemize}

In addition, to further reduce noise, the method of so-called SDAEs is utilized. The result
of the last hidden layer is then fed into a LSTM layer in order allowing the agent to
resolve the partial observability and discover latent patterns. For the model 5 stocks were
chosen (3 NASDAQ, 1 S\&P 500 stock index mini future, 1 HS300 stock index future). The
results applied on their test set show that both algorithms output significantly better
results in terms of annualized returns than both the original A3C and DQN algorithms as well
as simple buy and hold. Furthermore, their study shows that the SDAEs-LSTM A3C algorithm
performs significantly better than the SDAEs-LSTM DQN algorithm across all stocks.

\noindent\textit{Y. Li, W. Zheng and Z. Zheng, ``Deep Robust Reinforcement Learning for
Practical Algorithmic Trading,'' in IEEE Access, vol.~7, pp.~108014--108022, 2019,
doi: 10.1109/ACCESS.2019.2932789.}

\subsection{Related Work: Three Answers to the Generalization Problem}
\label{sec:related-work}

The three papers read for this report attack the same difficulty -- the gap between
in-sample fit and out-of-sample performance -- from three different sides, and all three
support the central observation of the 2026 update.

\paragraph{Representation.} Li, Zheng and Zheng (2019) change the input side of the agent:
stacked denoising autoencoders (SDAEs) clean the noisy market data and an LSTM layer
resolves the partial observability of the market state. On ten years of minute data their
extended DQN and A3C variants beat the plain algorithms and buy-and-hold in annualized
return. Their evaluation, however, still rests on a single 90/10 train--test split, so the
reported edge is an in-sample statement in the sense of our own comparison.

\paragraph{Environment.} Kuo et al.\ (2021) change the training environment instead of the
network. A LOB-GAN learns the ordering behaviour of the limit order book; its generator,
combined with a matching engine, forms an interactive ``Virtual Market'' in which the
portfolio agent is trained. Training in this simulated market improves out-of-sample
portfolio performance by about four percent over other generalization strategies. This is
the mirror image of our own result: a static environment built from historical prices is
itself part of the generalization problem.

\paragraph{Policy.} Niu, Li and Li (2022) keep the environment and diversify the policy set.
MetaTrader first imitates several expert strategies (momentum, mean reversion, hindsight) to
obtain a set of diverse base policies, and then trains a DQN meta-policy that recognises the
current market regime and switches to the most suitable base policy. On DJIA, HSI and CSI100
it reports clearly better risk-adjusted returns than DeepTrader and the classical
strategies.

Our own re-run (Chapter~\ref{sec:update-a2c} and the 20-seed extension in the accompanying
update report) sits between these three positions: the representation is unchanged (no
SDAE), the environment is still built from historical prices, and the policy set is a single
A2C agent. What was made robust instead is the evaluation itself -- every configuration was
repeated with twenty seeds. The outcome is sober and matches the papers' implicit lesson:
in-sample the agent wins clearly on all three stocks, but out-of-sample only Boeing keeps a
significant edge over buy-and-hold, Airbus ends exactly on buy-and-hold and Toyota turns
significantly worse. The in-sample margin does not predict the out-of-sample result. Read
against the three papers, this suggests that robustness has to be built into the
representation, the environment or the policy ensemble; a single-agent, single-environment
setup like ours should not be expected to generalise on its own.

\noindent\textit{C.-H. Kuo, C.-T. Chen, S.-J. Lin and S.-H. Huang, ``Improving Generalization
in Reinforcement Learning-Based Trading by Using a Generative Adversarial Market Model,'' in
IEEE Access, vol.~9, 2021, doi: 10.1109/ACCESS.2021.3068269.}

\noindent\textit{H. Niu, S. Li and J. Li, ``MetaTrader: A Reinforcement Learning Approach
Integrating Diverse Policies for Portfolio Optimization,'' in Proc.\ 31st ACM Int.\ Conf.\ on
Information and Knowledge Management (CIKM '22), Atlanta, GA, USA, 2022,
arXiv:2210.07771.}
